% Source: https://artofproblemsolving.com/community/c3h2372586_factors_and_lcm
% Stable identity: aops-topic-2372586-q01
A positive integer has exactly $3$ factors exactly when it is the square of a prime: if $n=p_1^{e_1}\cdots p_k^{e_k}$, it has $(e_1+1)\cdots(e_k+1)$ factors, and this product is $3$ only for $k=1$, $e_1=2$. The three-digit squares of primes are
\[11^2=121,\ 13^2=169,\ 17^2=289,\ 19^2=361,\ 23^2=529,\ 29^2=841,\ 31^2=961.\]
Next, $\overline{abcabc}=1000\cdot\overline{abc}+\overline{abc}=1001\cdot\overline{abc}=7\cdot11\cdot13\cdot p^2$.
\begin{itemize}
\item If $p\in\{17,19,23,29,31\}$: $7\cdot11\cdot13\cdot p^2$ has $2\cdot2\cdot2\cdot3=24$ factors.
\item If $p=11$: $7\cdot11^3\cdot13$ has $2\cdot4\cdot2=16$ factors; likewise $p=13$ gives $7\cdot11\cdot13^3$ with $16$ factors.
\end{itemize}
So $\overline{abcabc}$ has $\boxed{24\text{ factors, or }16\text{ factors when }\overline{abc}=121\text{ or }169}$.
